\documentclass[12pt,a4paper]{article}
\usepackage[utf8]{inputenc}
\usepackage[T1]{fontenc}
\usepackage{amsmath,amssymb,amsthm}
\usepackage{graphicx}
\usepackage{booktabs}
\usepackage{longtable}
\usepackage{array}
\usepackage{hyperref}
\usepackage{listings}
\usepackage{xcolor}
\usepackage[margin=1in]{geometry}
\usepackage{indentfirst}

\lstset{
    language=Python,
    basicstyle=\ttfamily\small,
    keywordstyle=\color{blue},
    commentstyle=\color{gray},
    stringstyle=\color{orange},
    breaklines=true,
    frame=single,
    showstringspaces=false
}

\title{Kuliah 09 - Pengenalan Machine Learning}
\author{Yeni Herdiyeni}
\date{}

\begin{document}
\maketitle
\tableofcontents
\newpage

\section{Kuliah 09 - Pengenalan Machine Learning}

\textbf{Topik:} Konsep Dasar, Tahapan, Evaluasi, dan Use Case  

\textbf{Pengampu:} Yeni Herdiyeni, Program Studi Kecerdasan Buatan, SSMI-IPB University  

\textbf{Tanggal:} 22 April 2026


\begin{quote}
\textbf{Catatan:} File asli berjudul \texttt{kuliah-09_pengenalan_machine_learning_slides.pdf}, namun di dalam slide tertulis "Kuliah 10". Dokumen ini mengikuti konteks materi pengenalan machine learning.\end{quote}


\hrule


\subsection{1. Capaian Pembelajaran}

Setelah mempelajari materi ini, mahasiswa diharapkan mampu:


\begin{enumerate}
\item Memahami tahapan utama dalam \emph{workflow} machine learning.
\end{enumerate}

\begin{enumerate}
\item Menjelaskan mengapa model perlu dievaluasi.
\end{enumerate}

\begin{enumerate}
\item Membedakan \textbf{underfitting}, \textbf{overfitting}, dan model yang seimbang.
\end{enumerate}

\begin{enumerate}
\item Memahami use case regresi sebagai contoh penerapan machine learning.
\end{enumerate}


\hrule


\subsection{2. Hubungan AI, Machine Learning, dan Data Science}

Tiga bidang ini sering tumpang tindih, tetapi memiliki cakupan yang berbeda:


\subsubsection{Artificial Intelligence (AI)}

\subsubsection{Machine Learning (ML)}

\subsubsection{Data Science}

\subsubsection{Tingkatan Data Analytics}

\begin{tabular}{|l|l|l|}
\hline
Tingkat & Pertanyaan Kunci & Contoh \\
\hline
\textbf{Descriptive} & Apa yang terjadi? & Produksi susu rata-rata 20 liter/hari \\
\hline
\textbf{Diagnostic} & Mengapa ini terjadi? & Mengapa hasil panen menurun? \\
\hline
\textbf{Predictive} & Apa yang akan terjadi? & Prediksi hasil panen musim depan \\
\hline
\textbf{Prescriptive} & Apa yang sebaiknya dilakukan? & Waktu tanam optimal \\
\hline
\end{tabular}


\hrule


\subsection{3. Mengapa Kita Perlu Machine Learning?}

\subsubsection{Use Case: Prediksi Produksi Tanaman}

\textbf{Input:}


\textbf{Output:} Hasil panen (ton/ha)


\paragraph{Pendekatan Tradisional}

\paragraph{Pendekatan Machine Learning}

\subsubsection{Definisi Machine Learning}

\begin{quote}
\textbf{Machine learning adalah pendekatan dalam kecerdasan buatan yang memungkinkan sistem belajar dari data untuk menemukan pola dan menghasilkan prediksi atau keputusan.}\end{quote}


Ide utama:


\hrule


\subsection{4. Konsep Learning dalam Machine Learning}

\subsubsection{Learning = Data + Model + Penyesuaian Parameter}

\textbf{Komponen utama:}


\\[
\hat{y} = f(X)
\\]


\textbf{Proses learning:}

\begin{enumerate}
\item Menghasilkan prediksi $\hat{y}$.
\end{enumerate}

\begin{enumerate}
\item Membandingkan prediksi dengan nilai sebenarnya $y$.
\end{enumerate}

\begin{enumerate}
\item Menghitung error.
\end{enumerate}

\begin{enumerate}
\item Menyesuaikan parameter model untuk meminimalkan error.
\end{enumerate}


\\[
\text{Error} = y - \hat{y}
\\]


\subsubsection{Perbandingan: Pemrograman Tradisional vs Machine Learning}

\begin{tabular}{|l|l|l|}
\hline
Aspek & Pemrograman Tradisional & Machine Learning \\
\hline
Input & Data + aturan eksplisit & Data + contoh output \\
\hline
Proses & Programmer menulis aturan & Sistem belajar pola sendiri \\
\hline
Output & Hasil dari aturan & Prediksi/keputusan dari model \\
\hline
Cocok untuk & Aturan yang jelas dan stabil & Aturan yang sulit ditulis manual \\
\hline
\end{tabular}


Contoh cocok untuk ML: prediksi harga rumah, deteksi penyakit, prediksi hasil panen, prediksi curah hujan.


\hrule


\subsection{5. Jenis-Jenis Learning}

\subsubsection{5.1 Supervised Learning (Pembelajaran Terbimbing)}

\textbf{Konsep:} Belajar dari data yang memiliki label/target.


\\[
(X, y) \rightarrow \text{Model} \rightarrow \hat{y}
\\]


\textbf{Tujuan:} Melakukan prediksi terhadap data baru.


\textbf{Contoh:}


\subsubsection{5.2 Unsupervised Learning (Pembelajaran Tak Terbimbing)}

\textbf{Konsep:} Belajar dari data tanpa label untuk menemukan pola tersembunyi.


\\[
X \rightarrow \text{Model} \rightarrow \text{Pola / Cluster}
\\]


\textbf{Tujuan:} Mengelompokkan atau menyederhanakan data.


\textbf{Contoh:}


\subsubsection{5.3 Reinforcement Learning (Pembelajaran Penguatan)}

\textbf{Konsep:} Belajar melalui interaksi dengan lingkungan, menggunakan aksi dan \emph{reward}.


\textbf{Komponen:}


\textbf{Contoh:}


\hrule


\subsection{6. Use Case Machine Learning dalam Kehidupan Nyata}

\begin{tabular}{|l|l|}
\hline
Domain & Aplikasi \\
\hline
\textbf{Pendidikan} & Prediksi mahasiswa berisiko terlambat lulus \\
\hline
\textbf{Kesehatan} & Prediksi kemungkinan penyakit dari hasil pemeriksaan \\
\hline
\textbf{Pertanian} & Prediksi hasil panen berdasarkan cuaca dan pupuk \\
\hline
\textbf{Kelautan} & Prediksi wilayah potensi penangkapan ikan dari data oseanografi \\
\hline
\textbf{Bisnis} & Rekomendasi produk dan prediksi churn pelanggan \\
\hline
\end{tabular}


\hrule


\subsection{7. Tahapan Machine Learning}

\subsubsection{7.1 Ringkasan Tahapan}

\begin{enumerate}
\item \textbf{Mengumpulkan data}
\end{enumerate}

\begin{enumerate}
\item \textbf{Pra-pemrosesan data}
\end{enumerate}

\begin{enumerate}
\item \textbf{Membagi data: training, validation, testing}
\end{enumerate}

\begin{enumerate}
\item \textbf{Melatih model}
\end{enumerate}

\begin{enumerate}
\item \textbf{Mengevaluasi model}
\end{enumerate}

\begin{enumerate}
\item \textbf{Memilih model terbaik}
\end{enumerate}

\begin{enumerate}
\item \textbf{Deployment dan monitoring}
\end{enumerate}


\subsubsection{7.2 Tahapan 1: Problem Formulation dan Data Preparation}

\textbf{Problem formulation:}


\textbf{Data preparation:}


\subsubsection{7.3 Tahapan 2: Training, Validation, dan Testing}

\begin{tabular}{|l|l|l|}
\hline
Subset & Fungsi & Proporsi Umum \\
\hline
\textbf{Training set} & Melatih model dan menyesuaikan parameter & 60-80% \\
\hline
\textbf{Validation set} & Memilih model atau hyperparameter terbaik & 10-20% \\
\hline
\textbf{Test set} & Menilai performa akhir pada data baru & 10-20% \\
\hline
\end{tabular}


\begin{quote}
\textbf{Intuisi:} training untuk belajar, validation untuk memilih, testing untuk menguji secara final.\end{quote}


\subsubsection{7.4 Tahapan 3: Deployment dan Monitoring}

\textbf{Deployment} berarti model digunakan pada lingkungan nyata.


Contoh deployment:


\textbf{Monitoring diperlukan karena:}


\hrule


\subsection{8. Mengapa Model Harus Dievaluasi?}

Model yang terlihat bagus pada data training belum tentu bagus pada data baru. Evaluasi model bertujuan untuk:


\begin{enumerate}
\item Membandingkan beberapa model.
\end{enumerate}

\begin{enumerate}
\item Mendeteksi underfitting atau overfitting.
\end{enumerate}

\begin{enumerate}
\item Memastikan model layak digunakan.
\end{enumerate}


\subsubsection{8.1 Evaluasi pada Regresi}

\begin{tabular}{|l|l|l|}
\hline
Metrik & Rumus & Interpretasi \\
\hline
\textbf{MAE} & $\frac{1}{n} \sum_{i=1}^{n} \lvert y_i - \hat{y}_i \rvert$ & Rata-rata kesalahan absolut \\
\hline
\textbf{MSE} & $\frac{1}{n} \sum_{i=1}^{n} (y_i - \hat{y}_i)^2$ & Rata-rata kesalahan kuadrat, memberi penalti lebih besar pada error besar \\
\hline
\textbf{RMSE} & $\sqrt{\frac{1}{n} \sum_{i=1}^{n} (y_i - \hat{y}_i)^2}$ & Error dalam satuan asli target \\
\hline
\textbf{$R^2$} & $1 - \frac{SS_{\text{res}}}{SS_{\text{tot}}}$ & Proporsi variansi data yang dijelaskan model \\
\hline
\end{tabular}


\textbf{Intuisi umum:}


\subsubsection{8.2 Evaluasi pada Klasifikasi}

\begin{tabular}{|l|l|}
\hline
Metrik & Definisi \\
\hline
\textbf{Akurasi} & Proporsi prediksi benar secara keseluruhan \\
\hline
\textbf{Precision} & Dari yang diprediksi positif, berapa yang benar-benar positif \\
\hline
\textbf{Recall} & Dari yang benar-benar positif, berapa yang berhasil ditemukan \\
\hline
\textbf{F1-score} & Rata-rata harmonik precision dan recall \\
\hline
\textbf{Confusion Matrix} & Tabel perbandingan prediksi vs aktual per kelas \\
\hline
\end{tabular}


Contoh: pada deteksi penyakit, \textbf{recall} sering penting agar kasus positif tidak banyak terlewat.


\subsubsection{8.3 Use Case Evaluasi di Dunia Nyata}

\begin{tabular}{|l|l|}
\hline
Kasus & Metrik yang Cocok \\
\hline
Prediksi hasil panen & MAE, RMSE \\
\hline
Deteksi tanaman sakit vs sehat & Precision, Recall, F1-score \\
\hline
Prediksi wilayah penangkapan ikan & Error prediksi, akurasi klasifikasi, atau metrik spasial \\
\hline
\end{tabular}


\hrule


\subsection{9. Underfitting, Overfitting, dan Model Seimbang}

\subsubsection{9.1 Underfitting}

\textbf{Ciri-ciri:}


\textbf{Contoh:} menggunakan model linear sederhana untuk data yang pola hubungannya sangat non-linear.


\subsubsection{9.2 Overfitting}

\textbf{Ciri-ciri:}


\textbf{Penyebab umum:}


\subsubsection{9.3 Balanced Model (Model Seimbang)}

\textbf{Ciri-ciri:}


\textbf{Tujuan utama ML:} menemukan model yang seimbang antara akurasi dan generalisasi.


\subsubsection{9.4 Tabel Perbandingan}

\begin{tabular}{|l|l|l|l|}
\hline
Kondisi & Kompleksitas & Train Error & Test Error \\
\hline
Underfitting & Rendah & Tinggi & Tinggi \\
\hline
Balanced & Sedang & Rendah & Rendah \\
\hline
Overfitting & Tinggi & Sangat rendah & Tinggi \\
\hline
\end{tabular}


\subsubsection{9.5 Cara Memilih Model}

\begin{enumerate}
\item Evaluasi dengan validation set.
\end{enumerate}

\begin{enumerate}
\item Bandingkan beberapa model dan parameter.
\end{enumerate}

\begin{enumerate}
\item Pilih model terbaik pada data validasi, bukan hanya training.
\end{enumerate}

\begin{enumerate}
\item Gunakan test set hanya untuk evaluasi akhir.
\end{enumerate}


\hrule


\subsection{10. Contoh Penerapan: Regresi untuk Prediksi Hasil Panen Padi}

\subsubsection{10.1 Tujuan Regresi}

Regresi bertujuan memprediksi nilai numerik atau kontinu.


\textbf{Contoh target regresi:}


\subsubsection{10.2 Use Case: Prediksi Hasil Panen Padi}

Misalkan kita ingin memprediksi hasil panen padi (ton/ha) berdasarkan fitur:


Masalah ini adalah \textbf{supervised learning jenis regresi} karena target yang diprediksi berupa nilai kontinu.


\subsubsection{10.3 Representasi Data}

Setiap baris data memiliki:


\subsubsection{10.4 Pembagian Data}

\begin{tabular}{|l|l|l|}
\hline
Subset & Proporsi & Fungsi \\
\hline
Training & 60% & Melatih model \\
\hline
Validation & 20% & Memilih model dan hyperparameter \\
\hline
Testing & 20% & Evaluasi performa akhir \\
\hline
\end{tabular}


\subsubsection{10.5 Model Regresi}

Model mencoba memetakan fitur ke hasil panen:


\\[
\hat{y} = f(X)
\\]


Contoh model linear:


\\[
\hat{y} = w_1 \cdot \text{curah hujan} + w_2 \cdot \text{suhu} + w_3 \cdot \text{kelembaban tanah} + w_4 \cdot \text{pupuk}
\\]


Model memiliki parameter $w$ yang dipelajari dari data training.


\subsubsection{10.6 Perhitungan Metrik (Contoh Numerik)}

\textbf{Data:}


\textbf{Error:}


\textbf{MAE:}


\\[
\begin{tabular}{|l|l|l|l|l|}
\hline
\text{MAE} = \frac{ & 0.5 & + & 2.0 & }{2} = 1.25 \\
\hline
\end{tabular}
\\]


\textbf{MSE:}


\\[
\text{MSE} = \frac{0.5^2 + 2.0^2}{2} = \frac{0.25 + 4}{2} = 2.125
\\]


\textbf{RMSE:}


\\[
\text{RMSE} = \sqrt{2.125} \approx 1.46
\\]


\textbf{MAPE:}


\\[
\text{MAPE} = \frac{1}{2} \left( \frac{0.5}{6.0} + \frac{2.0}{5.0} \right) \times 100\% \approx 24.2\%
\\]


\subsubsection{10.7 Interpretasi Metrik dalam Konteks Pertanian}

\begin{tabular}{|l|l|l|}
\hline
Metrik & Nilai & Makna \\
\hline
MAE & 1.25 ton/ha & Rata-rata kesalahan prediksi \\
\hline
RMSE & 1.46 ton/ha & Error dalam satuan asli, lebih mudah dipahami \\
\hline
MAPE & 24.2% & Persentase error, mudah dikomunikasikan ke petani \\
\hline
$R^2$ & - & Seberapa baik model mengikuti pola data \\
\hline
\end{tabular}


\subsubsection{10.8 Prediksi untuk Data Baru}

Setelah model dilatih dan dievaluasi, kita dapat memprediksi hasil panen untuk lahan baru:


\textbf{Contoh input baru:}


\\[
\hat{y} = f(X_{\text{baru}})
\\]


\textbf{Manfaat:}


\hrule


\subsection{11. Alur Berpikir pada Use Case Regresi}

\begin{enumerate}
\item Kumpulkan data historis.
\end{enumerate}

\begin{enumerate}
\item Bagi data menjadi training, validation, testing.
\end{enumerate}

\begin{enumerate}
\item Latih model regresi.
\end{enumerate}

\begin{enumerate}
\item Evaluasi error prediksi.
\end{enumerate}

\begin{enumerate}
\item Pilih model terbaik.
\end{enumerate}

\begin{enumerate}
\item Gunakan model untuk prediksi musim berikutnya.
\end{enumerate}


\hrule


\subsection{12. Ringkasan}

\begin{enumerate}
\item \textbf{Workflow ML} mencakup problem formulation, data collection, preprocessing, training, validation, testing, dan deployment.
\end{enumerate}

\begin{enumerate}
\item \textbf{Evaluasi model} penting untuk memastikan generalisasi, bukan hanya akurasi training.
\end{enumerate}

\begin{enumerate}
\item \textbf{Pemilihan model} harus mempertimbangkan underfitting, overfitting, dan keseimbangan.
\end{enumerate}

\begin{enumerate}
\item \textbf{Regresi} adalah contoh supervised learning untuk memprediksi nilai kontinu.
\end{enumerate}

\begin{enumerate}
\item Metrik seperti MAE, MSE, RMSE, dan $R^2$ membantu mengukur kualitas model regresi.
\end{enumerate}

\begin{enumerate}
\item Machine learning sangat berguna ketika pola data sulit dirumuskan dengan aturan manual.
\end{enumerate}



\end{document}
